\section{evnx backup / restore}

\begin{frame}{\cmd{evnx backup} \& \cmd{evnx restore}}
  \textbf{Purpose.} An encrypted copy of a \file{.env} you can store anywhere.

  \vspace{0.4em}
  \begin{columns}[T]
    \begin{column}{0.49\textwidth}
      \texttt{evnx backup [ENV] [OPTS]}
      \begin{tabular}{@{}l@{}}
        \flag{--output <FILE>} \\
        \flag{--key-file <PATH>} \\
        \flag{--keep <N>} \\
        \flag{--env-name <NAME>} \\
        \flag{--verify} \\
      \end{tabular}
    \end{column}
    \begin{column}{0.49\textwidth}
      \texttt{evnx restore <BACKUP> [OPTS]}
      \begin{tabular}{@{}l@{}}
        \flag{--output <FILE>} \\
        \flag{--password-file <PATH>} \\
        \flag{--inspect} \\
        \flag{--dry-run} \\
      \end{tabular}
    \end{column}
  \end{columns}

  \vspace{0.6em}
  \begin{block}{Each side also accepts the other's spelling}
    \cmd{backup} documents \flag{--key-file} and \cmd{restore} documents
    \flag{--password-file}, but \emph{both} accept \emph{both} — verified — and
    both honour \texttt{EVNX\_PASSWORD}. One key file always round-trips,
    including a binary one.
  \end{block}

  \vspace{0.3em}
  {\footnotesize Crypto: Argon2id over the passphrase $\rightarrow$ AES-256-GCM.
  Output is Base64.}
\end{frame}

\begin{frame}[fragile]{Backing up}
  \capture{53-backup.txt}
\end{frame}

\begin{frame}[fragile]{\flag{--verify}: decrypt it before trusting it}
  \capture{53b-backup-verify.txt}

  {\footnotesize Without \flag{--verify} the summary reports
  \texttt{Verified: no} — the backup was written but never read back. The advice
  ``test the restore before deleting the original'' is only as good as the test
  you actually run.}
\end{frame}

\begin{frame}[fragile]{\flag{--inspect}: what is in this backup?}
  \capture{54-restore-inspect.txt}

  {\footnotesize Key names only. Values are never printed, nothing is written,
  and no overwrite prompt appears. This is the right way to check a backup
  before committing to it.}
\end{frame}

\begin{frame}[fragile]{The round trip, verified}
  \capture{56-restore.txt}
  \vspace{-0.3em}
  \capture{57-restore-diff.txt}
\end{frame}

\begin{frame}{Exit codes are richer here — deliberately}
  \begin{tabular}{@{}cl@{}}
    \toprule
    \textbf{Code} & \textbf{Meaning} \\
    \midrule
    \exitzero & Success \\
    \exittwo  & Wrong password or corrupt backup \\
    \texttt{3} & Backup file not found \\
    \texttt{4} & Cancelled (informational) \\
    \texttt{5} & Validation fallback (informational) \\
    \texttt{6} & Integrity check failed \\
    \bottomrule
  \end{tabular}

  \vspace{0.8em}
  \textbf{Use cases}
  \begin{itemize}\footnotesize
    \item Nightly, with rotation: \texttt{evnx backup --env-name production --key-file /run/secrets/k --keep 7}
    \item CI restore with no file on disk: \texttt{EVNX\_PASSWORD="\$P" evnx restore prod.backup}
    \item Before a risky edit: \texttt{evnx backup \&\& \ldots}
  \end{itemize}

  \vspace{0.4em}
  \begin{alertblock}{There is deliberately no \flag{--force} on restore}
    Restoring over an existing file always prompts. The passphrase cannot be
    recovered — if you lose it, the backup is gone.
  \end{alertblock}
\end{frame}
